% =====================================================================
% iclr2027/theory.tex  --  body of \section{Theoretical Analysis}
%
% Usage in iclr2027_conference.tex:
%     \section{Theoretical Analysis}
%     \documentclass[10pt]{article}
\usepackage[margin=0.85in]{geometry}
\usepackage{amsmath,amssymb,amsthm}
\usepackage{graphicx}
\usepackage{caption}
\usepackage{kotex}
\captionsetup{font=small}

\newtheorem{proposition}{Proposition}

\newcommand{\supIPM}{\mathrm{supIPM}}
\newcommand{\EIPM}{\mathrm{EIPM}}
\newcommand{\hsupIPM}{\widehat{\supIPM}}
\newcommand{\Wone}{W_1}

\pagestyle{empty}
\begin{document}

\begin{center}\vspace{-0.5em}
{\large\bfseries Experiment}
\vspace{-0.5em}\end{center}


% Paste-ready theoretical-results section for the supIPM-FRL manuscript.
% Requires amsmath, amssymb, and amsthm, with lemma, theorem, and assumption
% environments defined in the preamble.

\section{Theoretical Guarantees}
\label{sec:theoretical-guarantees}

This section establishes two complementary properties of the proposed
criterion.  First, a small ReLU-supIPM controls the worst-case demographic
disparity of every downstream prediction head with bounded H\"older norm.
Second, the kernel estimator in Equation~(17) is uniformly consistent over a
possibly growing class of encoders.  The latter uniformity is essential: it
makes the result applicable to the data-dependent encoder returned by the
training procedure in Section~3.3.

Let
\[
  \tau(z):=\frac{z}{\sqrt{1+\lVert z\rVert_2^2}},
  \qquad \widetilde h:=\tau\circ h,
\]
so that \(\widetilde h(x)\in\mathbb B^d\).  Both the prediction head and the
ReLU discriminators below are applied to this same normalized
representation, and the representation dimension \(d\) is treated as fixed.
For an encoder \(h\), write
\begin{align}
  D_{\mathrm R}(h)
  &:=%
  \operatorname{supIPM}_{\mathcal V_{\mathrm R}}
  \bigl(\widetilde h(X);S\bigr) \notag\\
  &=\sup_{s\in\mathcal S}\sup_{v\in\mathcal V_{\mathrm R}}
  \Bigl|
    \mathbb E\!\left[v\!\left(\widetilde h(X)\right)\mid S=s\right]
    \notag\\[-1mm]
  &\hspace{42mm}
    -\mathbb E\!\left[v\!\left(\widetilde h(X)\right)\right]
  \Bigr| .
  \label{eq:population-relu-supipm}
\end{align}
Here and below, the conditional expectations use the pointwise version fixed
in Equation~(11).  In particular, for every Borel set
\(A\subseteq\mathbb B^d\),
\[
  P_{\widetilde h(X)\mid S=s}(A)
  :=\frac{1}{p_S(s)}
    \int_{\mathcal X}
    \mathbf 1\{\widetilde h(x)\in A\}p(x,s)\,d\nu(x).
\]
We denote the exact empirical criterion in Equation~(17),
viewed as a function of \(h\), by
\begin{align}
  \widehat D_{n,\gamma}(h)
  &:=%
  \sup_{s\in\mathcal S}\sup_{v\in\mathcal V_{\mathrm R}}
  \left|
    \sum_{i=1}^n
    \left\{\widehat w_\gamma(i;s)-\frac1n\right\}
    v\!\left(\widetilde h(X_i)\right)
  \right| .
  \label{eq:empirical-relu-supipm-theory}
\end{align}
For definiteness, a ratio with a zero kernel denominator is assigned the
value zero.  The event in Theorem~\ref{thm:uniform-supipm} has all kernel
denominators uniformly bounded away from zero, so this convention does not
affect any of the conclusions below.

For \(B>0\), let
\[
  \mathcal H_d^\beta(B)
  :=B\mathcal H_d^\beta
  =\left\{f:\mathbb B^d\to\mathbb R:
           \lVert f\rVert_{\mathcal H^\beta}\le B\right\}
\]
be the H\"older ball of radius \(B\).  For
\(\beta\ne(d+3)/2\), define
\[
  \alpha_{d,\beta}:=
  \begin{cases}
    \dfrac{2\beta}{d+3}, & 0<\beta<\dfrac{d+3}{2},\\[2mm]
    1, & \beta>\dfrac{d+3}{2}.
  \end{cases}
\]

\begin{lemma}[Worst-case fairness transfer]
\label{lem:worst-case-fairness-transfer}
Let \(h:\mathcal X\to\mathbb R^d\) be measurable and let
\(\beta>0\), \(\beta\ne(d+3)/2\).  Use the pointwise conditional-law
version fixed in Equation~(11).  There exists a constant
\(C_{d,\beta}>0\), depending only on \(d\) and \(\beta\), such that,
for every \(B>0\),
\begin{equation}
  \sup_{f\in\mathcal H_d^\beta(B)}
  \operatorname{GDP}_\infty(f\circ\widetilde h)
  \le
  C_{d,\beta}B\,D_{\mathrm R}(h)^{\alpha_{d,\beta}}.
  \label{eq:holder-head-fairness-bound}
\end{equation}
Consequently, a single bound on \(D_{\mathrm R}(h)\) holds
simultaneously for every \(s\in\mathcal S\) and every prediction head in
\(\mathcal H_d^\beta(B)\).  Moreover, if \(D_{\mathrm R}(h)=0\), then
\[
  P_{\widetilde h(X)\mid S=s}=P_{\widetilde h(X)}
  \qquad\text{for every }s\in\mathcal S,
\]
and hence perfect demographic parity holds for every bounded measurable
downstream head.
\end{lemma}

\begin{proof}
Fix \(s\in\mathcal S\) and set
\(P_s=P_{\widetilde h(X)\mid S=s}\) and
\(P=P_{\widetilde h(X)}\).  The ReLU--H\"older comparison inequality in
Section~2.2 gives, for every \(f\in\mathcal H_d^\beta(B)\),
\begin{align*}
  |P_sf-Pf|
  &\le B\operatorname{IPM}_{\mathcal H_d^\beta}(P_s,P)\\
  &\le C_{d,\beta}B
       \operatorname{IPM}_{\mathcal V_{\mathrm R}}(P_s,P)^%
       {\alpha_{d,\beta}}.
\end{align*}
Taking the supremum first over \(s\) and then over \(f\) proves
\eqref{eq:holder-head-fairness-bound}.  The final assertion follows because
ReLU-IPM separates probability measures on \(\mathbb B^d\).
\end{proof}

We next state conditions under which the sample criterion consistently
estimates \eqref{eq:population-relu-supipm}.  Let \(\mathcal E_n\) be a
deterministic encoder class and define the induced score class
\begin{equation}
  \mathcal G_n
  :=\left\{x\mapsto v\!\left(\widetilde h(x)\right):
      h\in\mathcal E_n,\ v\in\mathcal V_{\mathrm R}\right\}.
  \label{eq:induced-score-class}
\end{equation}
Every member of \(\mathcal G_n\) takes values in \([0,2]\).

\begin{assumption}[Sensitive-variable and kernel regularity]
\label{ass:kernel-regularity}
The observations \(\{(X_i,S_i,Y_i)\}_{i=1}^n\) are i.i.d.,
\(S\) is supported on \(\mathcal S=[a,b]\), and its density satisfies
\[
  0<c_S\le p_S(s)\le C_S<\infty,
  \qquad s\in\mathcal S.
\]
For \(g\in\bigcup_n\mathcal G_n\), let
\(m_g(s):=\mathbb E[g(X)\mid S=s]\), using the fixed pointwise version.
There is a constant \(L_S<\infty\), independent of \(n\), such that
\begin{equation}
  \sup_{g\in\mathcal G_n}|m_g(s)-m_g(t)|
  \le L_S|s-t|,
  \qquad s,t\in\mathcal S.
  \label{eq:uniform-conditional-lipschitz}
\end{equation}
The kernel \(K\) is bounded, nonnegative, symmetric, and of bounded
variation, with
\begin{align*}
  \int_{\mathbb R}K(u)\,du&=1,
  &\int_{\mathbb R}K(u)^2\,du&<\infty,\\
  &&\int_{\mathbb R}u^2K(u)\,du&<\infty.
\end{align*}
Write
\(K_\gamma(u,s)=\gamma^{-1}K((u-s)/\gamma)\), equivalently
\(K_\gamma(u-s)=\gamma^{-1}K((u-s)/\gamma)\).  The factor
\(\gamma^{-1}\) cancels from the normalized weights in Equation~(16).
Suppose that, for some \(\gamma_0>0\),
\begin{equation}
  \inf_{0<\gamma\le\gamma_0}\inf_{s\in\mathcal S}
  \int_a^b K_\gamma(u,s)\,du>0.
  \label{eq:one-sided-kernel-mass}
\end{equation}
\end{assumption}

Condition~\eqref{eq:uniform-conditional-lipschitz} is implied by the density
formulation in Equation~(11) if its smoothness assumption is strengthened to
the integrable-envelope condition
\begin{equation}
  \sum_{j=1}^2\int_{\mathcal X}
  \sup_{s\in\mathcal S}|\partial_s^j p(x,s)|\,d\nu(x)<\infty.
  \label{eq:density-envelope-condition}
\end{equation}
This strengthening is needed because pointwise boundedness in \(x\) alone
does not justify a uniform bound after integration over \(\mathcal X\).

\begin{assumption}[Complexity of the learned representations]
\label{ass:encoder-complexity}
The class \(\mathcal G_n\) is pointwise measurable.  There exist
\(A_n\ge e\) and \(\mathfrak v_n\ge1\) such that, for every finitely supported
probability measure \(Q\) and every \(0<\varepsilon\le1\),
\begin{equation}
  N\!\left(
    \varepsilon\lVert G_n\rVert_{Q,2},
    \mathcal G_n,L_2(Q)
  \right)
  \le \left(\frac{A_n}{\varepsilon}\right)^{\mathfrak v_n},
  \qquad G_n\equiv2.
  \label{eq:encoder-vc-entropy}
\end{equation}
The kernel-translate class
\(\mathcal K_\gamma:=\{t\mapsto K_\gamma(t,s):s\in\mathcal S\}\)
is pointwise measurable and is VC-type, uniformly in
\(0<\gamma\le\gamma_0\), with fixed constants \(A_K\ge e\) and
\(\mathfrak v_K\ge1\).  In particular, with envelope
\(K_\gamma^\star\equiv\lVert K\rVert_\infty/\gamma\),
\begin{equation}
  \sup_Q N\!\left(
    \varepsilon\lVert K_\gamma^\star\rVert_{Q,2},
    \mathcal K_\gamma,L_2(Q)
  \right)
  \le\left(\frac{A_K}{\varepsilon}\right)^{\mathfrak v_K}.
  \label{eq:kernel-vc-entropy}
\end{equation}
\end{assumption}

For \(0<\delta<1\), define
\begin{align}
  \Gamma_{n,\delta}
  &:=(\mathfrak v_n+\mathfrak v_K)
     \log\!\left(\frac{A_nA_Kn}{\gamma_n}\right)
     +\log\!\left(\frac{8}{\delta}\right),
  \label{eq:complexity-sequence}\\
  r_{n,\delta}
  &:=
       \gamma_n
       +\sqrt{\frac{\Gamma_{n,\delta}}{n\gamma_n}}
       +\frac{\Gamma_{n,\delta}}{n\gamma_n}
     .
  \label{eq:uniform-supipm-rate}
\end{align}

\begin{theorem}[Uniform consistency and a post-training fairness certificate]
\label{thm:uniform-supipm}
Suppose Assumptions~\ref{ass:kernel-regularity} and
\ref{ass:encoder-complexity} hold.  There exist constants \(C,c>0\),
independent of \(n\), \(\delta\), and \(h\), such that, whenever
  \(0<\gamma_n\le\gamma_0\wedge1\) and
\(\Gamma_{n,\delta}\le c n\gamma_n\), with probability at least
\(1-\delta\),
\begin{equation}
  \sup_{h\in\mathcal E_n}
  \left|\widehat D_{n,\gamma_n}(h)-D_{\mathrm R}(h)\right|
  \le C r_{n,\delta}.
  \label{eq:uniform-supipm-convergence}
\end{equation}
In particular, if \(\gamma_n\to0\) and there is a sequence
\(\delta_n\downarrow0\) such that
\[
  \frac{(\mathfrak v_n+\mathfrak v_K)\log(A_nA_Kn/\gamma_n)
        +\log(8/\delta_n)}{n\gamma_n}\to0,
\]
then the left-hand side of \eqref{eq:uniform-supipm-convergence} converges
to zero with probability at least \(1-\delta_n\), and hence is
\(o_{\mathbb P}(1)\).
Consequently, on the same event, for every \(\rho\ge C r_{n,\delta}\),
\begin{equation}
  \mathcal E_n(\rho-Cr_{n,\delta})
  \subseteq
  \widehat{\mathcal E}_{n,\gamma_n}(\rho)
  \subseteq
  \mathcal E_n(\rho+Cr_{n,\delta}),
  \label{eq:constraint-sandwich}
\end{equation}
where
\begin{align*}
  \mathcal E_n(\rho)
  &:=\{h\in\mathcal E_n:D_{\mathrm R}(h)\le\rho\},\\
  \widehat{\mathcal E}_{n,\gamma_n}(\rho)
  &:=\{h\in\mathcal E_n:
       \widehat D_{n,\gamma_n}(h)\le\rho\}.
\end{align*}
In particular, for any possibly data-dependent
\(\widehat h_n\in\mathcal E_n\), including an encoder returned by
Equation~(18) whenever its architecture and parameter constraints are
contained in \(\mathcal E_n\),
and every \(B>0\) and
\(\beta>0\), \(\beta\ne(d+3)/2\),
write
\[
  \Delta_{\beta,B}(h)
  :=\sup_{f\in\mathcal H_d^\beta(B)}
    \operatorname{GDP}_\infty(f\circ\tau\circ h).
\]
Then
\begin{align}
  \Delta_{\beta,B}(\widehat h_n)
  &\le C_{d,\beta}B\notag\\[-1mm]
  &\quad\times
  \left\{
    \widehat D_{n,\gamma_n}(\widehat h_n)+Cr_{n,\delta}
  \right\}^{\alpha_{d,\beta}}.
  \label{eq:post-training-fairness-certificate}
\end{align}
Thus, if
\(\widehat D_{n,\gamma_n}(\widehat h_n)\le\rho_n\), then, simultaneously
over all heads in \(\mathcal H_d^\beta(B)\),
\[
  \Delta_{\beta,B}(\widehat h_n)
  \le C_{d,\beta}B
  (\rho_n+Cr_{n,\delta})^{\alpha_{d,\beta}}.
\]
\end{theorem}

\begin{proof}
Let \(\mathbb P_n\) and \(P\) denote empirical and population averaging,
respectively.  For \(g\in\mathcal G_n\), introduce the population-smoothed
conditional mean
\begin{align*}
  m_{g,\gamma}(s)
  &:=\frac{\mathbb E[K_\gamma(S,s)g(X)]}
           {\mathbb E[K_\gamma(S,s)]},\\
  \widehat m_{g,\gamma}(s)
  &:=\frac{\mathbb P_n[K_\gamma(S,s)g(X)]}
           {\mathbb P_nK_\gamma(S,s)}.
\end{align*}
Set
\[
  \kappa_0:=
  \inf_{0<\gamma\le\gamma_0}\inf_{s\in\mathcal S}
  \int_a^bK_\gamma(u,s)\,du>0.
\]
For \(q_\gamma(s):=P K_\gamma(S,s)\), the density lower bound gives
\(q_\gamma(s)\ge c_S\kappa_0\).  Also,
\begin{align*}
  m_{g,\gamma}(s)-m_g(s)
  &=\int_a^b\frac{K_\gamma(u,s)p_S(u)}{q_\gamma(s)}\notag\\[-1mm]
  &\qquad\times\{m_g(u)-m_g(s)\}\,du.
\end{align*}
Consequently, the conditional Lipschitz assumption and the change of
variables \(t=(u-s)/\gamma\) imply, uniformly over
\((g,s)\in\mathcal G_n\times\mathcal S\),
\begin{equation}
  |m_{g,\gamma}(s)-m_g(s)|
  \le \frac{L_SC_S}{c_S\kappa_0}\,
       \gamma\int_{\mathbb R}|t|K(t)\,dt
  \le C\gamma.
  \label{eq:uniform-kernel-bias}
\end{equation}
Thus the bias bound also holds at the endpoints.

The elementary product-covering inequality applied to
\eqref{eq:encoder-vc-entropy} and \eqref{eq:kernel-vc-entropy} shows that
\[
  \bigl\{(x,u)\mapsto K_{\gamma_n}(u,s)g(x):
  (g,s)\in\mathcal G_n\times\mathcal S\bigr\}
\]
is VC-type with exponent of order
\(\mathfrak v_n+\mathfrak v_K\), envelope of order
\(\gamma_n^{-1}\), and second moment bounded by \(C\gamma_n^{-1}\).
A VC-type Bernstein maximal inequality, applied also to the kernel
denominator and to \(\mathcal G_n\), followed by a union bound, gives the
following statements simultaneously with probability at least
\(1-\delta\).  Uniformly over
\((g,s)\in\mathcal G_n\times\mathcal S\),
\begin{align*}
  \left|
    (\mathbb P_n-P)
    \{K_{\gamma_n}(S,s)g(X)\}
  \right|
  &\le C\sqrt{\frac{\Gamma_{n,\delta}}{n\gamma_n}}
       +C\frac{\Gamma_{n,\delta}}{n\gamma_n}.
\end{align*}
The same bound holds for
\(\sup_s|\mathbb P_nK_{\gamma_n}(S,s)-q_{\gamma_n}(s)|\), and
\begin{equation*}
  \sup_{g\in\mathcal G_n}|(\mathbb P_n-P)g|
  \le C\left\{
    \sqrt{\frac{\Gamma_{n,\delta}}{n}}
    +\frac{\Gamma_{n,\delta}}{n}
  \right\}.
\end{equation*}
The condition \(\Gamma_{n,\delta}\le cn\gamma_n\), for a sufficiently small
\(c\), therefore ensures
\[
  \inf_{s\in\mathcal S}\mathbb P_nK_{\gamma_n}(S,s)
  \ge \frac{c_S\kappa_0}{2}.
\]
Write
\(A_{g,s}:=(\mathbb P_n-P)\{K_{\gamma_n}(S,s)g(X)\}\) and
\(B_s:=(\mathbb P_n-P)K_{\gamma_n}(S,s)\).  Since
\(0\le m_{g,\gamma_n}(s)\le2\), the identity
\begin{align*}
  \widehat m_{g,\gamma_n}(s)-m_{g,\gamma_n}(s)
  =\frac{A_{g,s}-m_{g,\gamma_n}(s)B_s}
         {\mathbb P_nK_{\gamma_n}(S,s)}
\end{align*}
and \eqref{eq:uniform-kernel-bias} yield
\begin{equation*}
  \sup_{g\in\mathcal G_n}\sup_{s\in\mathcal S}
  \left|\widehat m_{g,\gamma_n}(s)-m_g(s)\right|
  \le C r_{n,\delta}.
\end{equation*}
For \(g_{h,v}(x):=v(\widetilde h(x))\), the inequality
\(\left|\sup_a|F(a)|-\sup_a|G(a)|\right|
\le\sup_a|F(a)-G(a)|\) now gives, uniformly in
\(h\in\mathcal E_n\),
\begin{align*}
  |\widehat D_{n,\gamma_n}(h)-D_{\mathrm R}(h)|
  &\le
  \sup_{g\in\mathcal G_n}\sup_{s\in\mathcal S}
  |\widehat m_{g,\gamma_n}(s)-m_g(s)|\\
  &\quad+
  \sup_{g\in\mathcal G_n}|(\mathbb P_n-P)g|
  \le Cr_{n,\delta},
\end{align*}
which proves \eqref{eq:uniform-supipm-convergence}.  Its convergence claim
follows by substituting \(\delta_n\) in the displayed bound.  The two set
inclusions are immediate from the same uniform event.  Finally,
\(D_{\mathrm R}(\widehat h_n)\le
\widehat D_{n,\gamma_n}(\widehat h_n)+Cr_{n,\delta}\); applying
Lemma~\ref{lem:worst-case-fairness-transfer} proves
\eqref{eq:post-training-fairness-certificate}.
\end{proof}

\paragraph{Rate and boundary effect.}
The term \(\gamma_n\) in \eqref{eq:uniform-supipm-rate} is the uniform bias
of the normalized local-constant estimator over the entire closed interval
\([a,b]\).  Thus the theorem genuinely covers every sensitive value in
\(\mathcal S\).  For a fixed-complexity encoder class, balancing this bias
with the stochastic term gives
  \(r_{n,\delta}=O((\log n/n)^{1/3})\), up to constants.  If, in
addition, \(K\) is supported on \([-1,1]\), the supremum is restricted
to an interior interval
\([a+\epsilon_n,b-\epsilon_n]\), with
\(\gamma_n<\epsilon_n\), and both \(p_S\) and the conditional score
functions have uniformly bounded first two derivatives, then a symmetric
second-order kernel gives
bias \(O(\gamma_n^2)\).  In that case \(\gamma_n\) in
\eqref{eq:uniform-supipm-rate} can be replaced by \(\gamma_n^2\), yielding
the familiar \(O((\log n/n)^{2/5})\) rate for a fixed-complexity
class.

\paragraph{Growing encoder classes.}
The quantities \(A_n\) and \(\mathfrak v_n\) may increase with \(n\).  Hence
Theorem~\ref{thm:uniform-supipm} allows a growing class of ReLU encoders
whenever
\[
  \gamma_n\to0,
  \qquad
  \frac{(\mathfrak v_n+\mathfrak v_K)
  \log(A_nA_Kn/\gamma_n)}{n\gamma_n}\to0.
\]
Here \(\mathfrak v_n\) may be taken to be any valid VC-subgraph,
pseudo-dimension, or covering-entropy upper bound for the composed class
\(\mathcal G_n\), including the entropy contributed by
\((\theta,\mu)\in\mathbb S^{d-1}\times[-1,1]\).  The radial map \(\tau\)
and each fixed ReLU discriminator are bounded and Lipschitz on their
respective domains, so an encoder-class covering bound extends to
\(\mathcal G_n\) after adding this finite-dimensional parameter entropy.

For example, suppose that
\(\sup_{x\in\mathcal X}\lVert x\rVert_\infty\le R_X<\infty\) and
\(\mathcal E_n\) has a fixed affine--ReLU architecture at each \(n\), with
depth \(L_n\), maximum width \(H_n\), at most \(W_n\) free weights and
biases, and all parameters bounded entrywise by \(M_n\).  A direct
parameter-covering argument gives
\begin{align}
  A_n^{\mathrm{net}}
  &:=C(R_X+1)(L_n+1)(W_n+d+1)\notag\\[-1mm]
  &\quad\times
  \{(M_n\vee1)(H_n+1)\}^{L_n+1}.
  \label{eq:network-entropy-constant}
\end{align}
Then
\begin{equation*}
  \log N(\varepsilon,\mathcal G_n,\lVert\cdot\rVert_\infty)
  \le C(W_n+d+1)
  \log\!\left(\frac{A_n^{\mathrm{net}}}{\varepsilon}\right).
\end{equation*}
Thus one may take \(\mathfrak v_n\lesssim W_n+d+1\), with
\(\log A_n\lesssim\log A_n^{\mathrm{net}}\).  In particular, a sufficient
growth condition is
\[
  \frac{(W_n+d+1)
        \log(A_n^{\mathrm{net}}n/\gamma_n)}{n\gamma_n}
  \to0.
\]
For a finite union of architectures, its log-cardinality is added to the
covering entropy.

\paragraph{Exact versus numerically approximated suprema.}
Theorem~\ref{thm:uniform-supipm} concerns the exact full-sample supremum
\(\widehat D_{n,\gamma_n}\) in Equation~(17); it does not assume that the
finite-start gradient scheme in Equations~(19)--(21) finds the global
maximum.  If a computed value \(\widetilde D_{n,\gamma_n}^{\rm alg}\) is
accompanied by a certified one-sided gap
\[
  0\le
  \widehat D_{n,\gamma_n}(\widehat h_n)
  -\widetilde D_{n,\gamma_n}^{\rm alg}(\widehat h_n)
  \le\zeta_n,
\]
then \(\widehat D_{n,\gamma_n}(\widehat h_n)\) in
\eqref{eq:post-training-fairness-certificate} may be replaced by
\(\widetilde D_{n,\gamma_n}^{\rm alg}(\widehat h_n)+\zeta_n\).  No encoder
optimization-error assumption is needed when the certificate is stated in
terms of the exact empirical criterion evaluated at the returned encoder.

\end{document}
% and, after \appendix:
%     % =====================================================================
% iclr2027/theory_appendix.tex  --  proofs for Section~\ref{sec:theory}.
% \input this file after \appendix in iclr2027_conference.tex.
% =====================================================================
\section{Proofs for Section~\ref{sec:theory}}
\label{app:proofs}

Throughout the appendix, $\mathbb{E}$ and $\mathbb{E}_n$ denote population and empirical averaging, $\mathbb{E}_n g:=n^{-1}\sum_{i=1}^n g(\mathbf{X}_i,S_i)$, and $C$ denotes a constant depending only on $c_S$, $C_S$, $\lVert M_1\rVert_{L_1(\nu)}$, $\lVert M_2\rVert_{L_1(\nu)}$ and $K$, whose value may change from line to line.  To lighten notation we write
\[
  D_n(h):=\operatorname{supIPM}_{\mathcal{V}_{\mathrm{R}},n}(\widetilde{\mathbf{Z}}_h;S),
  \qquad
  \widehat{D}_n(h):=\widehat{\operatorname{supIPM}}^{\gamma_n}_{\mathcal{V}_{\mathrm{R}},n}(\widetilde{\mathbf{Z}}_h;S),
\]
and $\gamma:=\gamma_n$.

\subsection{Proof of Theorem~\ref{thm:transfer}}

Fix $s\in\mathcal{S}_n$ and write $\mathbb{P}_s:=\mathbb{P}_{\widetilde{\mathbf{Z}}_h\mid S=s}$ and $\mathbb{P}_0:=\mathbb{P}_{\widetilde{\mathbf{Z}}_h}$, both probability measures on $\mathbb{B}^d$.  For $f\in\mathcal{H}^{\beta}_d(B)$, $f/B\in\mathcal{H}^{\beta}_d$, so Equation~\ref{eq:relu-holder} \citep[Theorem~2]{park2025relu}, whose proof requires $d\ge3$, yields
\[
  \left|\int f\,\mathrm{d}\mathbb{P}_s-\int f\,\mathrm{d}\mathbb{P}_0\right|
  \le B\operatorname{IPM}_{\mathcal{H}^{\beta}_d}(\mathbb{P}_s,\mathbb{P}_0)
  \le C_{d,\beta}\,B\,
     \operatorname{IPM}_{\mathcal{V}_{\mathrm{R}}}(\mathbb{P}_s,\mathbb{P}_0)^{\alpha_{d,\beta}}
  \le C_{d,\beta}\,B\,D_n(h)^{\alpha_{d,\beta}},
\]
where the last step uses $\operatorname{IPM}_{\mathcal{V}_{\mathrm{R}}}(\mathbb{P}_s,\mathbb{P}_0)\le D_n(h)$.  Taking the supremum over $s\in\mathcal{S}_n$ and then over $f$ proves Equation~\ref{eq:transfer}.  The bound for $v\in\mathcal{V}_{\mathrm{R}}$ is immediate from Equation~\ref{eq:trimmed-supipm}.  Finally, ReLU-IPM separates probability distributions on $\mathbb{B}^d$ \citep[Proposition~1]{park2025relu}, so $D_n(h)=0$ forces $\mathbb{P}_s=\mathbb{P}_0$ for every $s\in\mathcal{S}_n$.
\qed

\subsection{Proof of Theorem~\ref{thm:uniform}}
\label{app:proof-uniform}

Let $\mathcal{G}_n:=\{\mathbf{x}\mapsto v(\widetilde{h}(\mathbf{x})):h\in\mathcal{E}_n,\ v\in\mathcal{V}_{\mathrm{R}}\}$, so that $0\le g\le2$ for $g\in\mathcal{G}_n$.  For $g\in\mathcal{G}_n$ and $s\in\mathcal{S}$ set
\[
  a_g(s):=\int_{\mathcal{X}}g(\mathbf{x})p(\mathbf{x},s)\,\mathrm{d}\nu(\mathbf{x}),
  \qquad
  m_g(s):=\mathbb{E}[g(\mathbf{X})\mid S=s]=\frac{a_g(s)}{p_S(s)},
\]
which is the pointwise version in Equation~\ref{eq:pointwise-conditional}, and
\begin{align*}
  b_{g,\gamma}(s)&:=\mathbb{E}\{K_{\gamma}(S,s)g(\mathbf{X})\},&
  q_\gamma(s)&:=\mathbb{E}K_{\gamma}(S,s),\\
  m_{g,\gamma}(s)&:=\frac{b_{g,\gamma}(s)}{q_\gamma(s)},&
  \widehat{m}_{g,\gamma}(s)
  &:=\frac{\mathbb{E}_n\{K_{\gamma}(S,s)g(\mathbf{X})\}}{\mathbb{E}_nK_{\gamma}(S,s)}
   =\sum_{i=1}^n\widehat{w}_{\gamma}(i;s)\,g(\mathbf{X}_i).
\end{align*}
With this notation, $\widehat{D}_n(h)=\sup_{s,v}|\widehat{m}_{g,\gamma}(s)-\mathbb{E}_ng|$ and $D_n(h)=\sup_{s,v}|m_g(s)-\mathbb{E}g|$ for $g=v\circ\widetilde{h}$, the suprema running over $s\in\mathcal{S}_n$ and $v\in\mathcal{V}_{\mathrm{R}}$.

\paragraph{Step 0: entropy of the induced classes.}
The map $\tau$ is $1$-Lipschitz (its Jacobian is symmetric with eigenvalues in $(0,1]$), and for $\mathbf{z},\mathbf{z}'\in\mathbb{B}^d$,
$|v_{\theta,\mu}(\mathbf{z})-v_{\theta',\mu'}(\mathbf{z}')|\le\lVert\mathbf{z}-\mathbf{z}'\rVert_2+\lVert\theta-\theta'\rVert_2+|\mu-\mu'|$.  Hence
$\lVert v_{\theta,\mu}\circ\widetilde{h}-v_{\theta',\mu'}\circ\widetilde{h}'\rVert_\infty\le\lVert h-h'\rVert_\infty+\lVert\theta-\theta'\rVert_2+|\mu-\mu'|$, and covering $\mathbb{S}^{d-1}\times[-1,1]$ at scale $\eta/2$ gives
\[
  N(2\eta,\mathcal{G}_n,\lVert\cdot\rVert_\infty)
  \le N(\eta,\mathcal{E}_n,\lVert\cdot\rVert_\infty)\left(\frac{C_0}{\eta}\right)^{d+1}
  \le\left(\frac{A_n\vee C_0}{\eta}\right)^{\mathfrak{v}_n+d+1},
  \qquad 0<\eta\le1,
\]
for an absolute constant $C_0$.  Since $L_2(Q)$-distances are dominated by sup-norm distances and $\mathcal{G}_n$ has constant envelope $2$, $\mathcal{G}_n$ is VC-type with constants $(A_n\vee C_0,\ \mathfrak{v}_n+d+1)$.  By Assumption~\ref{as:kernel} and \citet[Lemma~22]{nolan1987u}, the class $\mathcal{K}_n:=\{s'\mapsto K_{\gamma}(s',s):s\in\mathcal{S}_n\}$ has envelope $\lVert K\rVert_\infty/\gamma$ and is VC-type with constants $(A_K,\mathfrak{v}_K)$ depending only on $K$.  By the product-covering inequality \citep[Corollary~A.1(i)]{chernozhukov2014gaussian},
\[
  \mathcal{F}_n:=\{(\mathbf{x},s')\mapsto K_{\gamma}(s',s)g(\mathbf{x}):
   g\in\mathcal{G}_n,\ s\in\mathcal{S}_n\}
\]
has envelope $2\lVert K\rVert_\infty/\gamma$ and is VC-type with exponent $\mathfrak{v}_n+d+1+\mathfrak{v}_K$ and constant $C(A_n\vee A_K\vee C_0)$.  Since $A_n\ge e$, $\gamma\le1$ and $\mathfrak{v}_n\ge1$, the resulting entropy factors satisfy
\begin{equation}
  (\mathfrak{v}_n+d+1+\mathfrak{v}_K)
  \log\frac{C(A_n\vee A_K\vee C_0)}{\gamma}
  \le C\,(\mathfrak{v}_n+d)\log\frac{A_n}{\gamma}
  =:C\,\Gamma_{n}^{0},
  \label{eq:entropy-factor}
\end{equation}
with $C$ depending only on $K$; note $\Gamma_n^0\ge2$.

\paragraph{Step 1: uniform second-order bias.}
Let $s\in\mathcal{S}_n$, $r_n:=\epsilon_n/\gamma_n\ge1$ and $I_n:=[-r_n,r_n]$, so that $s+\gamma t\in\mathcal{S}$ for every $t\in I_n$.  Extending $p(\mathbf{x},\cdot)$ and $p_S$ by zero outside $\mathcal{S}$, the substitution $s'=s+\gamma t$ gives
\begin{equation}
  b_{g,\gamma}(s)=\int_{\mathbb{R}}K(t)\,a_g(s+\gamma t)\,\mathrm{d}t,
  \qquad
  q_\gamma(s)=\int_{\mathbb{R}}K(t)\,p_S(s+\gamma t)\,\mathrm{d}t .
  \label{eq:full-mass}
\end{equation}
Since $0\le a_g\le2p_S\le2C_S$, the contribution of $t\notin I_n$ to either integral is at most $2C_ST_K(r_n)\le2C_S\gamma^2$ by Equation~\ref{eq:tail}; in particular $q_\gamma(s)\ge c_S\{1-T_K(r_n)\}\ge3c_S/4$ because $\gamma\le1/2$.  By dominated differentiation under Assumption~\ref{as:smooth-density}, $a_g$ and $p_S$ are twice continuously differentiable on $\mathcal{S}$ with
\begin{equation}
  \sup_{g\in\mathcal{G}_n}\lVert a_g''\rVert_\infty
  \le2\lVert M_2\rVert_{L_1(\nu)},
  \qquad
  \lVert p_S''\rVert_\infty\le\lVert M_2\rVert_{L_1(\nu)} .
  \label{eq:second-derivatives}
\end{equation}
On the symmetric interval $I_n$ a second-order Taylor expansion of $a_g(s+\gamma t)$ about $s$ gives
\[
  \int_{I_n}K(t)\,a_g(s+\gamma t)\,\mathrm{d}t-a_g(s)
  =-a_g(s)\,T_K(r_n)
   +\gamma\,a_g'(s)\int_{I_n}tK(t)\,\mathrm{d}t
   +\frac{\gamma^2}{2}\int_{I_n}t^2K(t)\,a_g''(\xi_t)\,\mathrm{d}t,
\]
where the middle term vanishes by symmetry.  Combining with the tail bound and Equation~\ref{eq:second-derivatives}, and arguing identically for $p_S$, we obtain uniformly over $g\in\mathcal{G}_n$ and $s\in\mathcal{S}_n$
\begin{equation}
  |b_{g,\gamma}(s)-a_g(s)|\le C\gamma^2,
  \qquad
  |q_\gamma(s)-p_S(s)|\le C\gamma^2 .
  \label{eq:bias}
\end{equation}
Because $0\le a_g\le2p_S$,
\begin{equation}
  |m_{g,\gamma}(s)-m_g(s)|
  \le\frac{|b_{g,\gamma}(s)-a_g(s)|}{q_\gamma(s)}
     +\frac{2\,|q_\gamma(s)-p_S(s)|}{q_\gamma(s)}
  \le C\gamma^2 .
  \label{eq:ratio-bias}
\end{equation}

\paragraph{Step 2: uniform stochastic error.}
For $f\in\mathcal{F}_n$, $\mathbb{E}f^2\le4\sup_{s}\mathbb{E}K_{\gamma}(S,s)^2\le4C_S\lVert K\rVert_\infty/\gamma$, so $\mathcal{F}_n$ has variance proxy $\sigma_n^2\asymp\gamma^{-1}$ and envelope of order $\gamma^{-1}$ (enlarging the envelope by a constant factor depending only on $K$ and $C_S$ if necessary, so that $\sigma_n$ does not exceed the $L_2$ norm of the envelope).  The maximal inequality of \citet[Corollary~5.1]{chernozhukov2014gaussian} and Equation~\ref{eq:entropy-factor} yield
\[
  \mathbb{E}\sup_{f\in\mathcal{F}_n}|(\mathbb{E}_n-\mathbb{E})f|
  \le C\left\{\sqrt{\frac{\Gamma_n^0}{n\gamma}}+\frac{\Gamma_n^0}{n\gamma}\right\}.
\]
Applying Bousquet's inequality \citep[Theorem~2.3]{bousquet2002bennett} to the uniformly bounded class $\mathcal{F}_n$ with $x=\log(3/\delta)$, and absorbing $\log3\le\Gamma_n^0$ into the constant, we obtain with probability at least $1-\delta/3$
\begin{equation}
  \sup_{\substack{g\in\mathcal{G}_n\\ s\in\mathcal{S}_n}}
  \left|(\mathbb{E}_n-\mathbb{E})\{K_{\gamma}(S,s)g(\mathbf{X})\}\right|
  \le C\left\{\sqrt{\frac{\Gamma_{n,\delta}}{n\gamma}}+\frac{\Gamma_{n,\delta}}{n\gamma}\right\}
  \le C\sqrt{\frac{\Gamma_{n,\delta}}{n\gamma}},
  \label{eq:num-stoch}
\end{equation}
where the last inequality uses $\Gamma_{n,\delta}\le c\,n\gamma$.  The same argument applied to $\mathcal{K}_n$ (no score factor) and to $\mathcal{G}_n$ (envelope $2$, variance proxy $4$, and $\gamma\le1$) gives, on an event of probability at least $1-\delta$ containing that of Equation~\ref{eq:num-stoch},
\begin{equation}
  \sup_{s\in\mathcal{S}_n}|(\mathbb{E}_n-\mathbb{E})K_{\gamma}(S,s)|
  \le C\sqrt{\frac{\Gamma_{n,\delta}}{n\gamma}},
  \qquad
  \sup_{g\in\mathcal{G}_n}|(\mathbb{E}_n-\mathbb{E})g|
  \le C\sqrt{\frac{\Gamma_{n,\delta}}{n\gamma}} .
  \label{eq:den-marg-stoch}
\end{equation}
By the lower bound $q_\gamma\ge3c_S/4$ from Step~1 and Equation~\ref{eq:den-marg-stoch}, choosing $c$ small enough that $C\sqrt{c}\le c_S/4$ ensures on this event
\begin{equation}
  \inf_{s\in\mathcal{S}_n}\mathbb{E}_nK_{\gamma}(S,s)\ \ge\ \frac{c_S}{2} .
  \label{eq:den-lower}
\end{equation}
With $A_{g,s}:=(\mathbb{E}_n-\mathbb{E})\{K_{\gamma}(S,s)g(\mathbf{X})\}$ and $B_s:=(\mathbb{E}_n-\mathbb{E})K_{\gamma}(S,s)$, the ratio identity
\[
  \widehat{m}_{g,\gamma}(s)-m_{g,\gamma}(s)
  =\frac{A_{g,s}-m_{g,\gamma}(s)\,B_s}{\mathbb{E}_nK_{\gamma}(S,s)},
  \qquad 0\le m_{g,\gamma}(s)\le2,
\]
together with Equations~\ref{eq:num-stoch}--\ref{eq:den-lower} and \ref{eq:ratio-bias} yields
\begin{equation}
  \sup_{\substack{g\in\mathcal{G}_n\\ s\in\mathcal{S}_n}}
  |\widehat{m}_{g,\gamma}(s)-m_g(s)|
  \le C\left\{\gamma^2+\sqrt{\frac{\Gamma_{n,\delta}}{n\gamma}}\right\}
  =C\,r_{n,\delta} .
  \label{eq:cond-mean-uniform}
\end{equation}

\paragraph{Step 3: conclusion.}
By the elementary inequality $|\sup_t|F(t)|-\sup_t|G(t)||\le\sup_t|F(t)-G(t)|$, for every $h\in\mathcal{E}_n$,
\[
  |\widehat{D}_n(h)-D_n(h)|
  \le\sup_{g,s}|\widehat{m}_{g,\gamma}(s)-m_g(s)|
     +\sup_{g}|(\mathbb{E}_n-\mathbb{E})g|
  \le C\,r_{n,\delta},
\]
by Equations~\ref{eq:cond-mean-uniform} and \ref{eq:den-marg-stoch}.  This proves Equation~\ref{eq:uniform}; the inclusions in Equation~\ref{eq:sandwich} follow from Equation~\ref{eq:uniform} on the same event.
\qed

\subsection{Proof of Corollary~\ref{cor:certificate}}

On the event of Theorem~\ref{thm:uniform}, $D_n(\widehat{h}_n)\le\widehat{D}_n(\widehat{h}_n)+Cr_{n,\delta}$, because Equation~\ref{eq:uniform} is uniform over $\mathcal{E}_n$ and hence applies to the data-dependent $\widehat{h}_n$.  Inserting this bound into Equation~\ref{eq:transfer} and using that $t\mapsto t^{\alpha_{d,\beta}}$ is increasing gives Equation~\ref{eq:certificate}.  For the rate, bounded $(A_n,\mathfrak{v}_n)$ and $\delta_n=n^{-c_1}$ give $\Gamma_{n,\delta_n}\asymp\log n$, so with $\gamma_n\asymp(\log n/n)^{1/5}$ both terms of $r_{n,\delta_n}$ are of order $(\log n/n)^{2/5}$; the conditions $\Gamma_{n,\delta_n}\le c\,n\gamma_n$ and $\gamma_n\le\epsilon_n\asymp1/\log n$ hold for all large $n$.
\qed

\subsection{Proof of Remark~\ref{rem:untrimmed}}
\label{app:untrimmed}

Let $\epsilon_n=0$ and $s\in\mathcal{S}$.  Put $T_s:=\{t\in\mathbb{R}:s+\gamma t\in\mathcal{S}\}$ and $\kappa(s):=\int_{T_s}K(t)\,\mathrm{d}t$.  The interval $T_s$ contains $[0,(b-a)/(2\gamma)]$ or $[-(b-a)/(2\gamma),0]$, so by symmetry $\kappa(s)\ge1/2-T_K\{(b-a)/(2\gamma)\}\ge1/4$ for all $n$ large enough that $T_K\{(b-a)/(2\gamma_n)\}\le1/4$ (for compactly supported $K$, as soon as $\gamma_n\le(b-a)/2$).  A first-order expansion with $\lVert a_g'\rVert_\infty\le2\lVert M_1\rVert_{L_1(\nu)}$ and $\lVert p_S'\rVert_\infty\le\lVert M_1\rVert_{L_1(\nu)}$ gives
\[
  b_{g,\gamma}(s)=\kappa(s)a_g(s)+R_b,\qquad
  q_\gamma(s)=\kappa(s)p_S(s)+R_q,
\]
with $|R_b|\vee|R_q|\le2\lVert M_1\rVert_{L_1(\nu)}\,\gamma\int|t|K(t)\,\mathrm{d}t$,
where $\int|t|K<\infty$ by Assumption~\ref{as:kernel}, and $q_\gamma(s)\ge c_S\kappa(s)\ge c_S/4$.  Hence
\[
  |m_{g,\gamma}(s)-m_g(s)|
  =\frac{|R_b\,p_S(s)-a_g(s)\,R_q|}{p_S(s)\,q_\gamma(s)}
  \le\frac{|R_b|+2|R_q|}{q_\gamma(s)}
  \le\frac{4\,(|R_b|+2|R_q|)}{c_S}
  \le C\gamma,
\]
uniformly over $g\in\mathcal{G}_n$ and $s\in\mathcal{S}$.  Steps~0,~2 and~3 of Appendix~\ref{app:proof-uniform} are unchanged (with $\mathcal{S}_n$ replaced by $\mathcal{S}$ and $c_S/2$ replaced by $c_S/8$ in Equation~\ref{eq:den-lower}), so the untrimmed estimator satisfies Equation~\ref{eq:uniform} with $\gamma_n^2$ replaced by $\gamma_n$ in $r_{n,\delta}$.
\qed

%
% Preamble additions (iclr2027_conference.tex already defines `assumption'):
%     \newtheorem{theorem}{Theorem}
%     \newtheorem{corollary}[theorem]{Corollary}
%     \theoremstyle{remark}\newtheorem{remark}{Remark}
% Bib entries to append to bibliography.bib: see theory_refs.bib.
% Notation follows Sections 2-3 of the main text (bold vectors, \mathbb{P},
% supIPM / \widehat{supIPM}^{\gamma}, GDP_\infty, \widehat{w}_\gamma(i;s),
% \mathcal{H}^\beta_d, Assumption~\ref{as:positive-density}).
% =====================================================================
\label{sec:theory}

We establish two properties of supIPM with the ReLU discriminator class.  First, a small population supIPM controls $\operatorname{GDP}_{\infty}$ of every prediction head of bounded H\"older norm (Theorem~\ref{thm:transfer}).  Second, the empirical supIPM in Equation~\ref{eq:empirical-supipm} converges to the population supIPM in Equation~\ref{eq:supipm} uniformly over the encoder class, so the empirical and population fairness constraints are asymptotically equivalent and the fairness level attained on the sample transfers to the population (Theorem~\ref{thm:uniform} and Corollary~\ref{cor:certificate}).  Proofs are given in Appendix~\ref{app:proofs}.

\subsection{Trimmed criteria}
\label{sec:theory-setup}

Write $\mathcal{S}=[a,b]$ and $\tau(\mathbf{z}):=\mathbf{z}/\sqrt{\lVert\mathbf{z}\rVert_2^2+1}$.  For an encoder $h:\mathcal{X}\to\mathbb{R}^d$ let $\widetilde{h}:=\tau\circ h$ and $\widetilde{\mathbf{Z}}_h:=\widetilde{h}(\mathbf{X})\in\mathbb{B}^d$, so that $\widetilde{\mathbf{Z}}_{h_\phi}$ is the normalized representation of Section~\ref{sec:supipm-criterion}.  As there, prediction heads $f:\mathbb{B}^d\to\mathbb{R}$ and discriminators are applied to $\widetilde{\mathbf{Z}}_h$, and conditional expectations given $S=s$ are the pointwise versions in Equation~\ref{eq:pointwise-conditional}.  Every $v\in\mathcal{V}_{\mathrm{R}}$ satisfies $0\le v\le2$ on $\mathbb{B}^d$.

The kernel weights $\widehat{w}_{\gamma}(i;s)$ in Equation~\ref{eq:supipm-estimators} are a local-constant smoother in $s$.  Within distance $\gamma$ of an endpoint of $\mathcal{S}$ such a smoother receives only part of the kernel mass and incurs a bias of first order in $\gamma$, whereas in the interior the bias is of second order.  We therefore state the finite-sample results on the expanding interior region
\begin{equation}
  \mathcal{S}_n:=[a+\epsilon_n,\,b-\epsilon_n],
  \qquad 0<\epsilon_n<\tfrac{b-a}{2},\quad \epsilon_n\to0,
  \label{eq:trimmed-domain}
\end{equation}
and write $\operatorname{GDP}_{\infty,n}$, $\operatorname{supIPM}_{\mathcal{V},n}$ and $\widehat{\operatorname{supIPM}}^{\gamma}_{\mathcal{V}_{\mathrm{R}},n}$ for Equations~\ref{eq:linfty-gdp}, \ref{eq:supipm} and \ref{eq:empirical-supipm} with the supremum over $s\in\mathcal{S}$ replaced by the supremum over $s\in\mathcal{S}_n$.  Explicitly,
\begin{align}
  \operatorname{GDP}_{\infty,n}(f\circ\widetilde{h})
  &:=\sup_{s\in\mathcal{S}_n}
     \left|\mathbb{E}[f(\widetilde{\mathbf{Z}}_h)\mid S=s]
           -\mathbb{E}[f(\widetilde{\mathbf{Z}}_h)]\right|,
  \label{eq:trimmed-gdp}\\
  \operatorname{supIPM}_{\mathcal{V}_{\mathrm{R}},n}(\widetilde{\mathbf{Z}}_h;S)
  &:=\sup_{s\in\mathcal{S}_n}
     \operatorname{IPM}_{\mathcal{V}_{\mathrm{R}}}
     \left(\mathbb{P}_{\widetilde{\mathbf{Z}}_h\mid S=s},
           \mathbb{P}_{\widetilde{\mathbf{Z}}_h}\right),
  \label{eq:trimmed-supipm}\\
  \widehat{\operatorname{supIPM}}^{\gamma}_{\mathcal{V}_{\mathrm{R}},n}
  (\widetilde{\mathbf{Z}}_h;S)
  &:=\sup_{s\in\mathcal{S}_n}
     \sup_{v\in\mathcal{V}_{\mathrm{R}}}
     \left|\sum_{i=1}^n
       \left(\widehat{w}_{\gamma}(i;s)-\frac1n\right)
       v\big(\widetilde{h}(\mathbf{X}_i)\big)\right|,
  \label{eq:trimmed-empirical-supipm}
\end{align}
with the convention $0/0:=0$ in $\widehat{w}_{\gamma}(i;s)$.  Since $\epsilon_n\to0$, every fixed $s\in(a,b)$ belongs to $\mathcal{S}_n$ for all large $n$; Remark~\ref{rem:untrimmed} describes what is lost when $\epsilon_n=0$.

\subsection{Assumptions}
\label{sec:theory-assumptions}

Assumption~\ref{as:positive-density} is in force throughout.  Let $p(\mathbf{x},s)$ be the joint density of $(\mathbf{X},S)$ with respect to $\nu\otimes\mathrm{Leb}$ fixed in Section~\ref{sec:supipm-criterion}, with $p_S(s)=\int_{\mathcal{X}}p(\mathbf{x},s)\,\mathrm{d}\nu(\mathbf{x})$ for every $s\in\mathcal{S}$.

\begin{assumption}[Smooth joint density]
\label{as:smooth-density}
$\mathcal{D}=\{(\mathbf{X}_i,S_i,Y_i)\}_{i=1}^n$ consists of i.i.d.\ copies of $(\mathbf{X},S,Y)$.  For $\nu$-almost every $\mathbf{x}$, the map $s\mapsto p(\mathbf{x},s)$ is twice continuously differentiable on $\mathcal{S}$, and there exist $M_1,M_2\in L_1(\nu)$ such that $\sup_{s\in\mathcal{S}}|\partial_s^{j}p(\mathbf{x},s)|\le M_j(\mathbf{x})$ for $j\in\{1,2\}$.
\end{assumption}

Assumption~\ref{as:smooth-density} is the counterpart of Assumption~3 of \citet{kong2025fair}, stated with integrable rather than pointwise envelopes; the envelopes are exactly what is needed to differentiate $s\mapsto\int g(\mathbf{x})p(\mathbf{x},s)\,\mathrm{d}\nu(\mathbf{x})$ under the integral sign uniformly over bounded $g$.  No smoothness in $\mathbf{x}$ is required.

\begin{assumption}[Kernel and tuning sequences]
\label{as:kernel}
$K_{\gamma}(s',s)=\gamma^{-1}K((s'-s)/\gamma)$ for a function $K:\mathbb{R}\to[0,\infty)$ that is bounded, symmetric, continuous and of bounded variation, with $\int K=1$ and $\int t^2K(t)\,\mathrm{d}t<\infty$.  Writing $T_K(r):=\int_{|t|>r}K(t)\,\mathrm{d}t$ for its tail mass, the bandwidth and trimming width satisfy $0<\gamma_n\le\min\{\epsilon_n,1/2\}$, $\gamma_n\to0$, and
\begin{equation}
  T_K(\epsilon_n/\gamma_n)\le\gamma_n^2 .
  \label{eq:tail}
\end{equation}
\end{assumption}

The factor $\gamma^{-1}$ in $K_{\gamma}$ cancels in the weights $\widehat{w}_{\gamma}(i;s)$ and is a normalization for the proofs.  The tail condition in Equation~\ref{eq:tail} says that, for $s\in\mathcal{S}_n$, the kernel mass falling outside $[s-\epsilon_n,s+\epsilon_n]\subseteq\mathcal{S}$ is of the order of the squared-bandwidth bias.  For compactly supported kernels such as the Epanechnikov and biweight kernels it holds trivially, since $T_K(\epsilon_n/\gamma_n)=0$ whenever $\gamma_n\le\epsilon_n$.  For the Gaussian kernel $T_K(r)\le e^{-r^2/2}$ for $r\ge1$, so Equation~\ref{eq:tail} holds whenever $\epsilon_n/\gamma_n\ge2\sqrt{\log(1/\gamma_n)}$; with $\epsilon_n\asymp1/\log n$ this is satisfied by every polynomially decaying bandwidth.  Bounded variation of $K$ makes the class $\{s'\mapsto K_{\gamma}(s',s):s\in\mathcal{S}_n\}$ VC-type with constants $(A_K,\mathfrak{v}_K)$ depending only on $K$ \citep[Lemma~22]{nolan1987u}.

\begin{assumption}[Encoder class]
\label{as:encoder}
$\mathcal{E}_n$ is a class of measurable encoders $h:\mathcal{X}\to\mathbb{R}^d$ such that, with $\lVert h-h'\rVert_\infty:=\sup_{\mathbf{x}\in\mathcal{X}}\lVert h(\mathbf{x})-h'(\mathbf{x})\rVert_2$, there exist $A_n\ge e$ and $\mathfrak{v}_n\ge1$ with
\begin{equation}
  N(\eta,\mathcal{E}_n,\lVert\cdot\rVert_\infty)
  \le\left(\frac{A_n}{\eta}\right)^{\mathfrak{v}_n},
  \qquad 0<\eta\le1 .
  \label{eq:encoder-entropy}
\end{equation}
\end{assumption}

Assumption~\ref{as:encoder} is the sup-norm entropy condition that also enters Theorem~5 of \citet{kong2025fair}.  It holds for the parameterized encoders $h_{\phi}$ of Algorithm~\ref{alg:supipm-frl} when $h_{\phi}$ is a ReLU network of depth $L$ and width $H$ with parameters bounded by a constant, with $\mathfrak{v}_n$ of the order of the number of parameters and $\log A_n$ of the order of $L\log H$ \citep[e.g.][Lemma~5]{schmidthieber2020nonparametric}.\footnote{To avoid measurability issues we assume, as usual, that all function classes appearing below are pointwise measurable \citep[Section~2.3.3]{vaart1996weak}.}  For $0<\delta<1$ define
\begin{equation}
  \Gamma_{n,\delta}
  :=(\mathfrak{v}_n+d)\log\frac{A_n}{\gamma_n}+\log\frac{1}{\delta},
  \qquad
  r_{n,\delta}
  :=\gamma_n^2+\sqrt{\frac{\Gamma_{n,\delta}}{n\gamma_n}} .
  \label{eq:rate}
\end{equation}

\subsection{Fairness transfer}
\label{sec:theory-transfer}

For $B>0$ let $\mathcal{H}^{\beta}_d(B):=\{f\in C^{k}(\mathbb{B}^d):\lVert f\rVert_{\mathcal{H}^{\beta}}\le B\}$ be the H\"older ball of radius $B$ in Equation~\ref{eq:holder-ball}, so that $\mathcal{H}^{\beta}_d=\mathcal{H}^{\beta}_d(1)$, and for $\beta\ne(d+3)/2$ let $\alpha_{d,\beta}:=2\beta/(d+3)$ if $\beta<(d+3)/2$ and $\alpha_{d,\beta}:=1$ if $\beta>(d+3)/2$ be the exponent in Equation~\ref{eq:relu-holder}.

\begin{theorem}[Fairness transfer]
\label{thm:transfer}
Let $d\ge3$,\footnote{The restriction $d\ge3$ is inherited from the ReLU approximation result underlying Equation~\ref{eq:relu-holder} \citep[Theorem~2]{park2025relu}.} and let $h:\mathcal{X}\to\mathbb{R}^d$ be measurable.  For every $\beta>0$ with $\beta\ne(d+3)/2$ there is a constant $C_{d,\beta}$, depending only on $d$ and $\beta$, such that for every $B>0$
\begin{equation}
  \sup_{f\in\mathcal{H}^{\beta}_d(B)}
  \operatorname{GDP}_{\infty,n}(f\circ\widetilde{h})
  \le C_{d,\beta}\,B\,
  \operatorname{supIPM}_{\mathcal{V}_{\mathrm{R}},n}(\widetilde{\mathbf{Z}}_h;S)^{\alpha_{d,\beta}} .
  \label{eq:transfer}
\end{equation}
Moreover, $\operatorname{GDP}_{\infty,n}(v\circ\widetilde{h})\le\operatorname{supIPM}_{\mathcal{V}_{\mathrm{R}},n}(\widetilde{\mathbf{Z}}_h;S)$ for every $v\in\mathcal{V}_{\mathrm{R}}$, and $\operatorname{supIPM}_{\mathcal{V}_{\mathrm{R}},n}(\widetilde{\mathbf{Z}}_h;S)=0$ implies $\mathbb{P}_{\widetilde{\mathbf{Z}}_h\mid S=s}=\mathbb{P}_{\widetilde{\mathbf{Z}}_h}$ for every $s\in\mathcal{S}_n$.
\end{theorem}

Theorem~\ref{thm:transfer} is the $L_{\infty}$ analogue of Theorems~1--2 of \citet{kong2025fair}: the bound holds at every $s\in\mathcal{S}_n$ rather than on average over $S$, and it covers every prediction head of bounded H\"older norm rather than the discriminators alone.  For heads smoother than $(d+3)/2$ the control is linear in supIPM.  With $\epsilon_n=0$ the same statement holds on all of $\mathcal{S}$ for the untrimmed criteria of Section~\ref{sec:supipm-criterion}.

\subsection{Uniform convergence and certificate}
\label{sec:theory-uniform}

\begin{theorem}[Uniform convergence]
\label{thm:uniform}
Let Assumptions~\ref{as:smooth-density}--\ref{as:encoder} hold.  There are constants $C,c>0$, depending only on $c_S$, $C_S$, $\lVert M_2\rVert_{L_1(\nu)}$ and $K$, such that whenever $\Gamma_{n,\delta}\le c\,n\gamma_n$, with probability at least $1-\delta$,
\begin{equation}
  \sup_{h\in\mathcal{E}_n}
  \left|
    \widehat{\operatorname{supIPM}}^{\gamma_n}_{\mathcal{V}_{\mathrm{R}},n}
    (\widetilde{\mathbf{Z}}_h;S)
    -\operatorname{supIPM}_{\mathcal{V}_{\mathrm{R}},n}
    (\widetilde{\mathbf{Z}}_h;S)
  \right|
  \le C\,r_{n,\delta} .
  \label{eq:uniform}
\end{equation}
Consequently, on the same event, for every fairness level $\rho\ge Cr_{n,\delta}$,
\begin{equation}
  \mathcal{E}_n(\rho-Cr_{n,\delta})
  \subseteq\widehat{\mathcal{E}}_n(\rho)
  \subseteq\mathcal{E}_n(\rho+Cr_{n,\delta}),
  \label{eq:sandwich}
\end{equation}
where $\mathcal{E}_n(t):=\{h\in\mathcal{E}_n:\operatorname{supIPM}_{\mathcal{V}_{\mathrm{R}},n}(\widetilde{\mathbf{Z}}_h;S)\le t\}$ and $\widehat{\mathcal{E}}_n(t):=\{h\in\mathcal{E}_n:\widehat{\operatorname{supIPM}}^{\gamma_n}_{\mathcal{V}_{\mathrm{R}},n}(\widetilde{\mathbf{Z}}_h;S)\le t\}$.
\end{theorem}

The rate $r_{n,\delta}$ has the usual kernel-smoothing form: a squared-bandwidth bias and a variance term of order $(n\gamma_n)^{-1/2}$ inflated by the entropy of the encoder class.  The representation dimension $d$ enters only through the $d+1$ parameters of the ReLU discriminator, so there is no curse of dimensionality in $\mathcal{Z}$, as in the discussion following Theorem~4 of \citet{kong2025fair}.  For a single fixed encoder ($\mathcal{E}_n=\{h\}$, $\mathfrak{v}_n=1$, $A_n=e$) Equation~\ref{eq:uniform} is the analogue of their Theorem~4, and Equation~\ref{eq:sandwich} of their Theorem~5: every encoder that satisfies the empirical constraint $\widehat{\operatorname{supIPM}}^{\gamma_n}_{\mathcal{V}_{\mathrm{R}},n}\le\rho$ satisfies the population constraint at level $\rho+Cr_{n,\delta}$.

\begin{corollary}[Certificate for the learned encoder]
\label{cor:certificate}
Let the assumptions of Theorems~\ref{thm:transfer} and \ref{thm:uniform} hold and let $\widehat{h}_n\in\mathcal{E}_n$ be any encoder, possibly depending on $\mathcal{D}$, such as the output $h_{\phi}$ of Algorithm~\ref{alg:supipm-frl}.  On the event of Theorem~\ref{thm:uniform}, simultaneously for all $B>0$ and all $\beta>0$ with $\beta\ne(d+3)/2$,
\begin{equation}
  \sup_{f\in\mathcal{H}^{\beta}_d(B)}
  \operatorname{GDP}_{\infty,n}(f\circ\tau\circ\widehat{h}_n)
  \le C_{d,\beta}\,B\,
  \Big\{
    \widehat{\operatorname{supIPM}}^{\gamma_n}_{\mathcal{V}_{\mathrm{R}},n}
    (\widetilde{\mathbf{Z}}_{\widehat{h}_n};S)
    +Cr_{n,\delta}
  \Big\}^{\alpha_{d,\beta}} .
  \label{eq:certificate}
\end{equation}
If, in addition, $A_n$ and $\mathfrak{v}_n$ are bounded, $\delta_n=n^{-c_1}$ for some $c_1>0$, $\gamma_n\asymp(\log n/n)^{1/5}$, and $\epsilon_n$ satisfies Assumption~\ref{as:kernel} (for instance $\epsilon_n\asymp1/\log n$, which works for both compactly supported and Gaussian kernels), then $r_{n,\delta_n}=O\{(\log n/n)^{2/5}\}$.
\end{corollary}

Equation~\ref{eq:certificate} is the operational content of the theory: the empirical supIPM of the learned encoder bounds the population worst-case disparity of every H\"older prediction head on $\mathcal{S}_n$, up to $Cr_{n,\delta}$.

\begin{remark}[Estimation without trimming]
\label{rem:untrimmed}
If $\epsilon_n=0$, so that the suprema run over all of $\mathcal{S}$ as in Equations~\ref{eq:supipm} and \ref{eq:empirical-supipm}, the proof of Theorem~\ref{thm:uniform} goes through with one change: for $s$ within $\gamma_n$ of an endpoint the kernel mass inside $\mathcal{S}$ is only partial, so the first-order Taylor term no longer cancels and the bias becomes $O(\gamma_n)$ instead of $O(\gamma_n^2)$ (Appendix~\ref{app:untrimmed}).  The untrimmed estimator therefore converges uniformly over $\mathcal{S}$ at rate $\gamma_n+\sqrt{\Gamma_{n,\delta}/(n\gamma_n)}$, that is $O\{(\log n/n)^{1/3}\}$ for $\gamma_n\asymp(\log n/n)^{1/3}$.  Trimming affects the rate but not consistency.
\end{remark}

\begin{remark}[Numerical supremum and held-out evaluation]
\label{rem:practical}
Theorem~\ref{thm:uniform} concerns the exact supremum in Equation~\ref{eq:trimmed-empirical-supipm}.  The value $\max_{m\in[M]}\Delta_m(\hat{s})$ returned by Algorithm~\ref{alg:supipm-frl} is a lower bound, and Equation~\ref{eq:certificate} applies to it after adding a verified optimization gap.  If, alternatively, Equation~\ref{eq:trimmed-empirical-supipm} is evaluated on an independent sample of size $m$ after training, the encoder class conditional on $\mathcal{D}$ is the singleton $\{\widehat{h}_n\}$, and Theorem~\ref{thm:uniform} applies with $n$ replaced by $m$ and $\Gamma_{m,\delta}=(d+1)\log(e/\gamma_m)+\log(1/\delta)$, irrespective of the architecture or the number of parameters of $\widehat{h}_n$.
\end{remark}
